\textbf{What was learned.} The method of \citet{Kiss2024} is reproducible from the paper alone: every
qualitative finding and every number that does not depend on the stopping point agrees. The one
quantitative gap is in how far the optimisation is run, which the paper does not report. A reproduction
should therefore report the optimiser budget and the stopping rule as carefully as the model.

\textbf{Threats to validity.} (i) The authors' optimiser options, $\epsilon$, steady-state rule and
seeds are unknown; we fixed our choices in the M2 protocol before seeing results and tested the most
influential ones ($\epsilon$, optimiser, budget, grid, initial state). (ii) H1 is supported by two
independent budget-limited optimisers, but \texttt{fmincon}'s interior-point steps were not reproduced
exactly; without MATLAB it remains a hypothesis. (iii) The Fig.~8 values of the paper were read from the
axes. (iv) The extension uses one system, one model class and simulated white sensor noise. Noisy states
in the regressors could bias least squares (errors-in-variables, \citealp{Soderstrom2007}); E4 shows that this
bias is small here and that the noise error is mostly variance, amplified by $\sqrt2/T_s$ in the target. With
a fixed absolute noise level, designs that spread the states more also have a higher effective SNR; the
noise-free case (\mFOfullinf$\times$) isolates the coverage effect. (v) The surrogate design used short runs (20 iterations); longer runs on a fixed
surrogate helped little, and a second identification round appeared more useful (unequal budgets, not tested statistically). (vi) The converged
optima depend on the machine (\S3), so only statistics over seeds, not single converged values, should be
compared.

\textbf{Unsuccessful attempts} (kept in the package): $P=5$ violated the steady-state rule for the
optimised signal; a first gradient check with a too small step reported false errors; PyTorch could
not be installed, so all gradients were derived by hand. An early test of the extension pipeline on seed 0 used
state-relative SNR, which penalises designs that excite the system more. After seeing that test we changed to
fixed absolute sensor noise, because a real sensor's noise does not depend on the input; this was a choice
made after an early result, and the noise-free case, which does not depend on it, shows the same ranking.

\textbf{Research use.} The extension suggests a practical recipe for robotics and UAV identification,
where the dynamics are unknown and test time is expensive: one short random-phase experiment, a local
model, then one or two space-filling redesigns on the identified model.
